\begin{frame}{Wie können \alert{IoT und LoRaWAN} die Forschung der Stakeholder unterstützen?}{Die Grundidee}
\begin{description}[Messung der Wasserchemie]
\item[Messung der Wasserchemie] \begin{itemize}
    \item Platzierung von LoRaWAN-Gateways rund um den Hartbeespoort-Damm, ausgestattet mit einem
    \begin{itemize}
        \item GPS-Empfänger, um deren exakte Position zu bestimmen, und
        \item Solarzellen und Akkumulatoren, um die Gateways autonom zu betreiben, und
        \item Integration in \emph{The Things Network}.
    \end{itemize}
    \item Entwicklung von schwimmenden LoRaWAN-Nodes, ausgestattet mit
    \begin{itemize}
        \item GPS-Empfänger, um dessen exakte Position zu bestimmen, und
        \item Solarzellen und Akkumulatoren, um den Node autonom zu betreiben
    \end{itemize}
    zur Messung von
    \begin{itemize}
        \item Wasserqualität
        \item Nährstoffgehalt im Wasser
        \item \ldots
    \end{itemize}
    \item Integration in das \emph{The Things Network}.
    \item Abholen der Daten per MQTT und Speicherung dieser in MariaDB.
    \item Entwicklung einer Weboberfläche für Management und Datenerfassung.
\end{itemize}
\end{description}
\end{frame}
